% GENERATED FILE. Edit canonical lessons, then run npm run generate:books.
% canonical-source-hash: fnv1a64:b2708e951c24b365
% canonical-lessons: MR-C145-gal, MR-C145-jibh, MR-C145-hanuvati, MR-C145-kopar, MR-C145-managat

\chapter{Cheek, Tongue, Chin, Elbow, Wrist}
\label{ch:mr-cheek-tongue-chin-elbow-wrist}
\hlchaptermodality{145}

\begin{chapteropening}
\textbf{By the end of this chapter:} \emph{I can say a cheek, the tongue, the chin, an elbow and a wrist.}

Complete the last lesson of chapter 145: I can say a cheek, the tongue, the chin, an elbow and a wrist.
\end{chapteropening}

\section[\texorpdfstring{\mr{गाल} (gāl)}{गाल (gāl)}]{\mr{गाल} (gāl) — a cheek}

\label{lesson:MR-C145-gal}



\begin{quote}
Before the new one: say the Marathi for a mosquito, then the Marathi for the forehead.
\end{quote}

\subsection*{You'll want to know: \mr{गाल}}
\textbf{\mr{गाल}} — \emph{gāl} — \textquotedblleft{}a cheek\textquotedblright{}.

\begin{cousinweb}
\textbf{\mr{गाल}} (\emph{gāl}) is a tadbhava of Sanskrit \emph{galla}, \textquotedblleft{}cheek\textquotedblright{} (CDIAL 4089). CDIAL notes the view that the word is from Dravidian.
\end{cousinweb}

\begin{grammarlens}[title={What you've built}]
One more word for this chapter.
\end{grammarlens}

\subsection*{Your turn}
\begin{itemize}
\raggedright
  \item \emph{Say it:} \emph{gāl}
  \item \emph{Say it:} \emph{gāl}, once more
  \item \emph{Say it:} say \emph{gāl}
  \item \emph{Recall:} say the Marathi for a mosquito, then the Marathi for the forehead, then say \emph{gāl} again
\end{itemize}

\begin{tcolorbox}[breakable,colback=teal!4,colframe=teal!35!black,title={Before you move on}]
What does \mr{गाल} mean? (\textquotedblleft{}A cheek\textquotedblright{}.) Say it once more.
\end{tcolorbox}
\section[\texorpdfstring{\mr{जीभ} (jībh)}{जीभ (jībh)}]{\mr{जीभ} (jībh) — the tongue}

\label{lesson:MR-C145-jibh}



\begin{quote}
Before the new one: say the Marathi for the forehead, then the Marathi for a cheek.
\end{quote}

\subsection*{You'll want to know: \mr{जीभ}}
\textbf{\mr{जीभ}} — \emph{jībh} — \textquotedblleft{}the tongue\textquotedblright{}.

\begin{cousinweb}
\textbf{\mr{जीभ}} (\emph{jībh}) is a tadbhava, a word inherited from Sanskrit and reshaped by sound change. It comes from \emph{jihvā}, \textquotedblleft{}tongue\textquotedblright{} (CDIAL 5228).
\end{cousinweb}

\begin{grammarlens}[title={What you've built}]
One more word for this chapter.
\end{grammarlens}

\subsection*{Your turn}
\begin{itemize}
\raggedright
  \item \emph{Say it:} \emph{jībh}
  \item \emph{Say it:} \emph{jībh}, once more
  \item \emph{Say it:} say \emph{jībh}
  \item \emph{Recall:} say the Marathi for the forehead, then the Marathi for a cheek, then say \emph{jībh} again
\end{itemize}

\begin{tcolorbox}[breakable,colback=teal!4,colframe=teal!35!black,title={Before you move on}]
What does \mr{जीभ} mean? (\textquotedblleft{}The tongue\textquotedblright{}.) Say it once more.
\end{tcolorbox}
\section[\texorpdfstring{\mr{हनुवटी} (hanuvaṭī)}{हनुवटी (hanuvaṭī)}]{\mr{हनुवटी} (hanuvaṭī) — the chin}

\label{lesson:MR-C145-hanuvati}



\begin{quote}
Before the new one: say the Marathi for a cheek, then the Marathi for the tongue.
\end{quote}

\subsection*{You'll want to know: \mr{हनुवटी}}
\textbf{\mr{हनुवटी}} — \emph{hanuvaṭī} — \textquotedblleft{}the chin\textquotedblright{}.

\begin{cousinweb}
\textbf{\mr{हनुवटी}} (\emph{hanuvaṭī}) is a tadbhava, a word inherited from Sanskrit and reshaped by sound change. It comes from \emph{hanu}, \textquotedblleft{}jaw\textquotedblright{} (CDIAL 13966). CDIAL lists Marathi \emph{haṇū} and \emph{haṇvaṭī}, \textquotedblleft{}chin\textquotedblright{}.
\end{cousinweb}

\begin{grammarlens}[title={What you've built}]
One more word for this chapter.
\end{grammarlens}

\subsection*{Your turn}
\begin{itemize}
\raggedright
  \item \emph{Say it:} \emph{hanuvaṭī}
  \item \emph{Say it:} \emph{hanuvaṭī}, once more
  \item \emph{Say it:} say \emph{hanuvaṭī}
  \item \emph{Recall:} say the Marathi for a cheek, then the Marathi for the tongue, then say \emph{hanuvaṭī} again
\end{itemize}

\begin{tcolorbox}[breakable,colback=teal!4,colframe=teal!35!black,title={Before you move on}]
What does \mr{हनुवटी} mean? (\textquotedblleft{}The chin\textquotedblright{}.) Say it once more.
\end{tcolorbox}
\section[\texorpdfstring{\mr{कोपर} (kopar)}{कोपर (kopar)}]{\mr{कोपर} (kopar) — an elbow}

\label{lesson:MR-C145-kopar}



\begin{quote}
Before the new one: say the Marathi for the tongue, then the Marathi for the chin.
\end{quote}

\subsection*{You'll want to know: \mr{कोपर}}
\textbf{\mr{कोपर}} — \emph{kopar} — \textquotedblleft{}an elbow\textquotedblright{}.

\begin{cousinweb}
\textbf{\mr{कोपर}} (\emph{kopar}) is a tadbhava, a word inherited from Sanskrit and reshaped by sound change. It comes from \emph{kūrpara}, \textquotedblleft{}elbow\textquotedblright{} (CDIAL 3413).
\end{cousinweb}

\begin{grammarlens}[title={What you've built}]
One more word for this chapter.
\end{grammarlens}

\subsection*{Your turn}
\begin{itemize}
\raggedright
  \item \emph{Say it:} \emph{kopar}
  \item \emph{Say it:} \emph{kopar}, once more
  \item \emph{Say it:} say \emph{kopar}
  \item \emph{Recall:} say the Marathi for the tongue, then the Marathi for the chin, then say \emph{kopar} again
\end{itemize}

\begin{tcolorbox}[breakable,colback=teal!4,colframe=teal!35!black,title={Before you move on}]
What does \mr{कोपर} mean? (\textquotedblleft{}An elbow\textquotedblright{}.) Say it once more.
\end{tcolorbox}
\section[\texorpdfstring{\mr{मनगट} (managaṭ)}{मनगट (managaṭ)}]{\mr{मनगट} (managaṭ) — a wrist}

\label{lesson:MR-C145-managat}



\begin{quote}
Before the new one: say the Marathi for the chin, then the Marathi for an elbow.
\end{quote}

\subsection*{You'll want to know: \mr{मनगट}}
\textbf{\mr{मनगट}} — \emph{managaṭ} — \textquotedblleft{}a wrist\textquotedblright{}.

\begin{cousinweb}
\textbf{\mr{मनगट}} (\emph{managaṭ}) is a tadbhava, a word inherited and reshaped by sound change. It comes from \emph{*maṇigranthi}, \textquotedblleft{}wrist\textquotedblright{}, a form CDIAL reconstructs from its descendants (CDIAL 9734). It is built on \emph{granthi}, \textquotedblleft{}knot, joint\textquotedblright{}.
\end{cousinweb}

\begin{grammarlens}[title={What you've built}]
One more word for this chapter.
\end{grammarlens}

\subsection*{Your turn}
\begin{itemize}
\raggedright
  \item \emph{Say it:} \emph{managaṭ}
  \item \emph{Say it:} \emph{managaṭ}, once more
  \item \emph{Say it:} say \emph{managaṭ}
  \item \emph{Recall:} say the Marathi for the chin, then the Marathi for an elbow, then say \emph{managaṭ} again
\end{itemize}

\begin{tcolorbox}[breakable,colback=teal!4,colframe=teal!35!black,title={Before you move on}]
What does \mr{मनगट} mean? (\textquotedblleft{}A wrist\textquotedblright{}.) Say it once more.
\end{tcolorbox}
